\begin{thebibliography}{99}
\bibitem[1]{published-1} S Aaronson and A Wigderson. Algebrization: A New Barrier in Complexity Theory. ACM Trans. Comput. Theory, 1(1), February 2009.
\bibitem[2]{published-2} E. Allender, M. Cheraghchi, D. Myrisiotis, H. Tirumala, and I. Volkovich. One-Way Functions and a Conditional Variant of MKTP. In 41st IARCS Annual Conference on Foundations of Software Technology and Theoretical Computer Science, FSTTCS 2021, December 15-17, 2021, Virtual Conference, volume 213 of LIPIcs, pages 7:1–7:19. Schloss Dagstuhl - Leibniz-Zentrum für Informatik, 2021.
\bibitem[3]{published-3} B. Applebaum, B. Barak, and D. Xiao. On Basing Lower-Bounds for Learning on Worst-Case Assumptions. In Proceedings of the 49th Annual IEEE Symposium on Foundations of Computer Science, FOCS’08, pages 211–220, 2008.
\bibitem[4]{published-4} T. Baker, J. Gill, and R. Solovay. Relativizations of the P =?N P Question. SIAM Journal on Computing, 4(4):431–442, 1975.
\bibitem[5]{published-5} Manuel Blum and Sampath Kannan. Designing Programs that Check Their Work. J. ACM, 42(1):269–291, 1995. doi:10.1145/200836.200880.
\bibitem[6]{published-6} Manuel Blum and Silvio Micali. How to Generate Cryptographically Strong Sequences of Pseudo-Random Bits. SIAM J. Comput., 13(4):850–864, 1984. doi:10.1137/0213053.
\bibitem[7]{published-7} A. Bogdanov and L. Trevisan. Average-Case Complexity. Foundations and Trends in Theoretical Computer Science, 2(1):1–106, 2006.
\bibitem[8]{published-8} G. Brassard. Relativized Cryptography. IEEE Transactions on Information Theory, 29(6):877– 894, 1983.
\bibitem[9]{published-9} Harry Buhrman, Lance Fortnow, and Aduri Pavan. Some Results on Derandomization. Theory Comput. Syst., 38(2):211–227, 2005. doi:10.1007/s00224-004-1194-y.
\bibitem[10]{published-10} M. Carmosino, R. Impagliazzo, V. Kabanets, and A. Kolokolova. Learning Algorithms from Natural Proofs. In Proceedings of the 31st Conference on Computational Complexity, CCC’16. Schloss Dagstuhl–Leibniz-Zentrum fuer Informatik, 2016.
\bibitem[11]{published-11} R. Gennaro and L. Trevisan. Lower bounds on the efficiency of generic cryptographic constructions. In Proceedings 41st Annual Symposium on Foundations of Computer Science, pages 305–313, 2000.
\bibitem[12]{published-12} O. Goldreich. Foundations of Cryptography: Volume 1. Cambridge University Press, New York, NY, USA, 2006.
\bibitem[13]{published-13} O. Goldreich, S. Goldwasser, and S. Micali. How to Construct Random Functions. J. ACM, 33(4):792–807, August 1986.
\bibitem[14]{published-14} J. Håstad. Computational Limitations of Small-depth Circuits. MIT Press, Cambridge, MA, USA, 1987.
\bibitem[15]{published-15} J. Håstad, R. Impagliazzo, L. A. Levin, and M. Luby. A Pseudorandom Generator from Any One-way Function. SIAM J. Comput., 28(4):1364–1396, March 1999.
\bibitem[16]{published-16} S. Hirahara. Non-Black-Box Worst-Case to Average-Case Reductions within NP. In 59th IEEE Annual Symposium on Foundations of Computer Science, FOCS 2018, Paris, France, October 7-9, 2018, pages 247–258, 2018.
\bibitem[17]{published-17} S Hirahara. Characterizing Average-Case Complexity of PH by Worst-Case Meta-Complexity. In IEEE 61st Annual Symposium on Foundations of Computer Science (FOCS 2020), pages 50–60, 2020.
\bibitem[18]{published-18} S. Hirahara. Non-Disjoint Promise Problems from Meta-Computational View of Pseudorandom Generator Constructions. In 35th Computational Complexity Conference (CCC 2020), volume 169 of LIPIcs, pages 20:1–20:47, Dagstuhl, Germany, 2020.
\bibitem[19]{published-19} S. Hirahara. Average-Case Hardness of NP from Exponential Worst-Case Hardness Assump- tions. In 53rd Annual ACM Symposium on Theory of Computing (STOC 2021), 2021.
\bibitem[20]{published-20} S. Hirahara and M. Nanashima. On Worst-Case Learning in Relativized Heuristica. In 62nd Annual IEEE Symposium on Foundations of Computer Science (FOCS 2021), 2021.
\bibitem[21]{published-21} S. Hirahara and R. Santhanam. On the Average-Case Complexity of MCSP and Its Variants. In Ryan O’Donnell, editor, 32nd Computational Complexity Conference, CCC 2017, July 6-9, 2017, Riga, Latvia, volume 79 of LIPIcs, pages 7:1–7:20. Schloss Dagstuhl - Leibniz-Zentrum für Informatik, 2017.
\bibitem[22]{published-22} S. Hirahara and R. Santhanam. Errorless versus Error-prone Average-case Complexity. In The 13th Innovations in Theoretical Computer Science (ITCS 2022), 2022.
\bibitem[23]{published-23} R. Ilango, H. Ren, and R. Santhanam. Hardness on any Samplable Distribution Suffices: New Characterizations of One-Way Functions by Meta-Complexity. Electron. Colloquium Comput. Complex., page 82, 2021.
\bibitem[24]{published-24} R. Impagliazzo. A personal view of average-case complexity. In Proceedings of IEEE Tenth Annual Conference on Structure in Complexity Theory, pages 134–147, 1995.
\bibitem[25]{published-25} R. Impagliazzo. Relativized Separations of Worst-Case and Average-Case Complexities for NP. In 2011 IEEE 26th Annual Conference on Computational Complexity, pages 104–114, 2011.
\bibitem[26]{published-26} R. Impagliazzo and L. Levin. No better ways to generate hard NP instances than picking uniformly at random. In Proceedings of the 31st Annual Symposium on Foundations of Computer Science, FOCS’90, pages 812–821, 1990.
\bibitem[27]{published-27} R. Impagliazzo and S. Rudich. Limits on the Provable Consequences of One-Way Permutations. In Proceedings of the Twenty-First Annual ACM Symposium on Theory of Computing, STOC ’89, pages 44–61, New York, NY, USA, 1989. ACM.
\bibitem[28]{published-28} K. Ko. On the Complexity of Learning Minimum Time-Bounded Turing Machines. SIAM Journal on Computing, 20(5):962–986, 1991.
\bibitem[29]{published-29} L. A Levin. Average Case Complete Problems. SIAM J. Comput., 15(1):285–286, February 1986.
\bibitem[30]{published-30} Y. Liu and R. Pass. On One-way Functions and Kolmogorov Complexity. In IEEE 61st Annual Symposium on Foundations of Computer Science (FOCS 2020), pages 1243–1254, 2020. doi:10.1109/FOCS46700.2020.00118.
\bibitem[31]{published-31} Y. Liu and R. Pass. A Note on One-way Functions and Sparse Languages. Electron. Colloquium Comput. Complex., page 92, 2021.
\bibitem[32]{published-32} Y. Liu and R. Pass. On the Possibility of Basing Cryptography on EXP $\neq$ BPP. In Advances in Cryptology - CRYPTO 2021 - 41st Annual International Cryptology Conference, CRYPTO 2021, Virtual Event, August 16-20, 2021, Proceedings, Part I, volume 12825 of Lecture Notes in Computer Science, pages 11–40. Springer, 2021.
\bibitem[33]{published-33} M. Nanashima. A Theory of Heuristic Learnability. In Proceedings of the 34th Conference on Learning Theory, COLT’21. PMLR, August 2021.
\bibitem[34]{published-34} M. Nanashima. On Basing Auxiliary-Input Cryptography on NP-Hardness via Nonadaptive Black-Box Reductions. In James R. Lee, editor, 12th Innovations in Theoretical Computer Science Conference (ITCS 2021), volume 185 of LIPIcs, pages 29:1–29:15, Dagstuhl, Germany, 2021. Schloss Dagstuhl–Leibniz-Zentrum für Informatik.
\bibitem[35]{published-35} I. Oliveira and R. Santhanam. Conspiracies between Learning Algorithms, Circuit Lower Bounds, and Pseudorandomness. In Proceedings of the 32nd Computational Complexity Conference, CCC’17, Dagstuhl, DEU, 2017. Schloss Dagstuhl–Leibniz-Zentrum fuer Informatik.
\bibitem[36]{published-36} R. Ostrovsky. One-way functions, hard on average problems, and statistical zero-knowledge proofs. In Proceedings of the Sixth Annual Structure in Complexity Theory Conference, pages 133–138, 1991.
\bibitem[37]{published-37} R. Ostrovsky and A. Wigderson. One-way functions are essential for non-trivial zero-knowledge. In Proceedings of the 2nd Israel Symposium on Theory and Computing Systems, ISTCS’93, pages 3–17, June 1993.
\bibitem[38]{published-38} A. A. Razborov and S. Rudich. Natural Proofs. Journal of Computer and System Sciences, 55(1):24–35, 1997.
\bibitem[39]{published-39} H. Ren and R. Santhanam. A Relativization Perspective on Meta-Complexity. In The 39th International Symposium on Theoretical Aspects of Computer Science (STACS 2022), 2022.
\bibitem[40]{published-40} L. Valiant. A Theory of the Learnable. Commun. ACM, 27(11):1134–1142, 1984. doi: 10.1145/1968.1972.
\bibitem[41]{published-41} Thomas Watson. Relativized Worlds without Worst-Case to Average-Case Reductions for NP. ACM Trans. Comput. Theory, 4(3), September 2012.
\bibitem[42]{published-42} H Wee. Finding Pessiland. In Proceedings of the Third Conference on Theory of Cryptography, TCC’06, pages 429–442, Berlin, Heidelberg, 2006. Springer-Verlag.
\bibitem[43]{published-43} D. Xiao. On basing ZK $\neq$ BPP on the hardness of PAC learning. In Proceedings of the 24th Conference on Computational Complexity, CCC’09, pages 304–315, 2009.
\bibitem[44]{published-44} A. Yao. Theory and Application of Trapdoor Functions. In Proceedings of the 23rd Annual Symposium on Foundations of Computer Science, FOCS’82, pages 80–91, November 1982.
\bibitem[45]{published-45} M. Zimand. Efficient Privatization of Random Bits. In IN WORKSHOP ON RANDOMIZED ALGORITHMS, 1997.
\end{thebibliography}
